
\documentclass{article}
\usepackage{colortbl}
\usepackage{makecell}
\usepackage{multirow}
\usepackage{supertabular}

\begin{document}

\newcounter{utterance}

\centering \large Interaction Transcript for game `codenames', experiment `abstract', episode 7 with reflexion{-}llama{-}3.1{-}8b{-}instruct.
\vspace{24pt}

{ \footnotesize  \setcounter{utterance}{1}
\setlength{\tabcolsep}{0pt}
\begin{supertabular}{c@{$\;$}|p{.15\linewidth}@{}p{.15\linewidth}p{.15\linewidth}p{.15\linewidth}p{.15\linewidth}p{.15\linewidth}}
    \# & $\;$A & \multicolumn{4}{c}{Game Master} & $\;\:$B\\
    \hline

    \theutterance \stepcounter{utterance}  
    & & & \multicolumn{2}{p{0.3\linewidth}}{
        \cellcolor[rgb]{0.9,0.9,0.9}{
            \makecell[{{p{\linewidth}}}]{
                \texttt{\tiny{[GM$|$GM]}}
                {'hidden': {'team': ['faithlessness', 'romanticism'], 'innocent': ['renown', 'antihuman', 'refutation', 'treachery'], 'opponent': [], 'assassin': ['insubordination']}, 'revealed': {'team': {'team': [], 'innocent': [], 'opponent': [], 'assassin': []}, 'opponent': {'team': [], 'innocent': [], 'opponent': [], 'assassin': []}}}
            }
        }
    }
    & & \\ \\

    \theutterance \stepcounter{utterance}  
    & & \multicolumn{4}{p{0.6\linewidth}}{
        \cellcolor[rgb]{0.9,0.9,0.9}{
            \makecell[{{p{\linewidth}}}]{
                \texttt{\tiny{[P1$\langle$GM]}}
                \texttt{You have attempted this same task before and did not succeed. These are the plans you wrote after those attempts.} \\
\texttt{Trial \#1: New plan: } \\
\\ 
\texttt{In my previous attempts, I focused on providing a single{-}word clue that relates to multiple words from my team's list. However, I failed to consider the potential impact of my clue on the opponent's words. In the last attempt, I provided a clue that was too broad, leading to the opponent guessing an innocent word (antipathy) and ending my teammate's turn.} \\
\\ 
\texttt{Next time, I will focus on providing a clue that is more specific and targeted towards my team's words, while avoiding any potential connections to the opponent's words. I will also consider the remaining words on the lists and try to find a clue that relates to multiple words from my team's list, but not to any of the opponent's words.} \\
\\ 
\texttt{To achieve this, I will analyze the remaining words on the lists and look for common themes, prefixes, or suffixes that can be used as a clue. I will also try to find a clue that is not too obvious, but still provides a clear connection to the target words.} \\
\\ 
\texttt{For example, if I notice that multiple words from my team's list have a similar prefix or suffix, I will use that as a clue. If I see a common theme or concept among the target words, I will try to find a word that relates to that theme.} \\
\\ 
\texttt{By being more strategic and targeted in my clue selection, I hope to avoid providing a clue that leads to the opponent guessing an innocent word and ending my teammate's turn.} \\
\\ 
\texttt{Trial \#2: New plan: In my previous attempts, I focused on providing a clue that relates to multiple words from my team's list, but I failed to consider the potential impact of my clue on the opponent's words. I also didn't analyze the remaining words on the lists to find a clue that is more specific and targeted towards my team's words.} \\
\\ 
\texttt{Next time, I will analyze the remaining words on the lists and look for common themes, prefixes, or suffixes that can be used as a clue. I will also try to find a clue that is not too obvious, but still provides a clear connection to the target words.} \\
\\ 
\texttt{To achieve this, I will:} \\
\\ 
\texttt{1. Analyze the prefixes and suffixes of the target words from my team's list.} \\
\texttt{2. Look for common themes or concepts among the target words.} \\
\texttt{3. Find a word that relates to the theme or concept, but is not too obvious.} \\
\texttt{4. Make sure the clue does not relate to any of the opponent's words.} \\
\\ 
\texttt{For example, if I notice that multiple words from my team's list have a similar prefix or suffix, I will use that as a clue. If I see a common theme or concept among the target words, I will try to find a word that relates to that theme.} \\
\\ 
\texttt{By being more strategic and targeted in my clue selection, I hope to provide a clue that is more effective in helping my teammate guess the target words, while avoiding any potential connections to the opponent's words.} \\
\\ 
\\ 
\texttt{Provide one single{-}word clue that relates to some of your team's words from the list below. You can choose to how many and to which words your clue relates to.} \\
\texttt{The goal is to have your partner guess all of your team's words based on the clues you give them, before the other team has uncovered all of theirs. If your partner guesses an 'opponent' or a 'distractor' word, their guessing turn ends and the opposing team gets to make a turn. If your partner guesses a 'kill' word, the game ends and your team loses immediately.} \\
\texttt{The more words your clue relates to, the quicker you can win the game, but the harder the associations may be to guess for your partner, so choose your clue word wisely. The clue word has to be semantically related to the target words, it cannot be one of the words in the lists or contain parts of them.} \\
\\ 
\texttt{Always give your single{-}word clue and your comma{-}separated list of related target words in the following format and make your answers as short as possible, never include any other text than is required in this form:} \\
\\ 
\texttt{CLUE: <WORD>} \\
\texttt{TARGETS: <TARGETS>} \\
\\ 
\texttt{Your team words are: faithlessness, romanticism.} \\
\texttt{Your opponent's team words are: .} \\
\texttt{Distractor words are: renown, antihuman, refutation, treachery.} \\
\texttt{The kill words are: insubordination.} \\
            }
        }
    }
    & \\ \\

    \theutterance \stepcounter{utterance}  
    & \multicolumn{4}{p{0.6\linewidth}}{
        \cellcolor[rgb]{0.9,0.9,0.9}{
            \makecell[{{p{\linewidth}}}]{
                \texttt{\tiny{[P1$\rangle$GM]}}
                \texttt{CLUE: Ideal} \\
\texttt{TARGETS: romanticism, faithlessness} \\
            }
        }
    }
    & & \\ \\

    \theutterance \stepcounter{utterance}  
    & & & \multicolumn{2}{p{0.3\linewidth}}{
        \cellcolor[rgb]{0.9,0.9,0.9}{
            \makecell[{{p{\linewidth}}}]{
                \texttt{\tiny{[GM$|$GM]}}
                \texttt{ideal} \\
            }
        }
    }
    & & \\ \\

    \theutterance \stepcounter{utterance}  
    & & & \multicolumn{2}{p{0.3\linewidth}}{
        \cellcolor[rgb]{0.9,0.9,0.9}{
            \makecell[{{p{\linewidth}}}]{
                \texttt{\tiny{[GM$|$GM]}}
                ['romanticism', 'faithlessness']
            }
        }
    }
    & & \\ \\

    \theutterance \stepcounter{utterance}  
    & & & \multicolumn{2}{p{0.3\linewidth}}{
        \cellcolor[rgb]{0.9,0.9,0.9}{
            \makecell[{{p{\linewidth}}}]{
                \texttt{\tiny{[GM$|$GM]}}
                {'word': 'romanticism', 'assignment': 'team'}
            }
        }
    }
    & & \\ \\

    \theutterance \stepcounter{utterance}  
    & & & \multicolumn{2}{p{0.3\linewidth}}{
        \cellcolor[rgb]{0.9,0.9,0.9}{
            \makecell[{{p{\linewidth}}}]{
                \texttt{\tiny{[GM$|$GM]}}
                {'word': 'faithlessness', 'assignment': 'team'}
            }
        }
    }
    & & \\ \\

    \theutterance \stepcounter{utterance}  
    & & & \multicolumn{2}{p{0.3\linewidth}}{
        \cellcolor[rgb]{0.9,0.9,0.9}{
            \makecell[{{p{\linewidth}}}]{
                \texttt{\tiny{[GM$|$GM]}}
                \texttt{opponent has won} \\
            }
        }
    }
    & & \\ \\

    \theutterance \stepcounter{utterance}  
    & \multicolumn{4}{p{0.6\linewidth}}{
        \cellcolor[rgb]{0.95,0.95,0.95}{
            \makecell[{{p{\linewidth}}}]{
                \texttt{\tiny{[P2$\rightarrow$P1]}}
                \texttt{Post{-}episode Reflexion (not sent to GM):} \\
\texttt{New plan: In my previous attempts, I focused on providing a clue that relates to multiple words from my team's list, but I failed to consider the potential impact of my clue on the opponent's words. I also didn't analyze the remaining words on the lists to find a clue that is more specific and targeted towards my team's words.} \\
\\ 
\texttt{Next time, I will analyze the remaining words on the lists and look for common themes, prefixes, or suffixes that can be used as a clue. I will also try to find a clue that is not too obvious, but still provides a clear connection to the target words.} \\
\\ 
\texttt{To achieve this, I will:} \\
\\ 
\texttt{1. Analyze the prefixes and suffixes of the target words from my team's list.} \\
\texttt{2. Look for common themes or concepts among the target words.} \\
\texttt{3. Find a word that relates to the theme or concept, but is not too obvious.} \\
\texttt{4. Make sure the clue does not relate to any of the opponent's words.} \\
\\ 
\texttt{Additionally, I will also consider the following:} \\
\\ 
\texttt{5. If I notice that multiple words from my team's list have a similar prefix or suffix, I will use that as a clue.} \\
\texttt{6. If I see a common theme or concept among the target words, I will try to find a word that relates to that theme.} \\
\texttt{7. I will also try to find a clue that is more specific and targeted towards my team's words, while avoiding any potential connections to the opponent's words.} \\
\\ 
\texttt{By being more strategic and targeted in my clue selection, I hope to provide a clue that is more effective in helping my teammate guess the target words, while avoiding any potential connections to the opponent's words.} \\
            }
        }
    }
    & & \\ \\

\end{supertabular}
}

\end{document}
